Unskilled and Unaware of It: How Difficulties in Recognizing One's Own
Incompetence Lead to Inflated Self-Assessments

Several lines of research are consistent with the notion that
incompetent individuals lack the metacognitive skills necessary for
accurate self-assessment. Work on the nature of expertise, for
instance, has revealed that novices possess poorer metacognitive
skills than do experts. In physics, novices are less accurate than
experts in judging the difficulty of physics problems (Chi et al.,
1982). In chess, novices are less calibrated than experts about how
many times they need to see a given chessboard position before
they are able to reproduce it correctly (Chi, 1978). In tennis,
novices are less likely than experts to successfully gauge whether
specific play attempts were successful (McPherson & Thomas,
1989).
These findings suggest that unaccomplished individuals do not
possess the degree of metacognitive skills necessary for accurate
self-assessment that their more accomplished counterparts possess.
However, none of this research has examined whether metacog-
nitive deficiencies translate into inflated self-assessments or
whether the relatively incompetent (novices) are systematically
more miscalibrated about their ability than are the competent
(experts).


A language model can produce a coherent response without providing a reliable guarantee that the response is factually correct,
within scope, or aligned with the decisions made by other agents [Who is in charge? Coordination and governance in agentic
software engineering Paolo Ciancarini]

Any well-calibrated language model assigns
non-zero probability to previously unseen combinations of words and facts. Plausible but unsupported output is therefore not merely
an occasional implementation fault; it is a possibility that must be managed by the surrounding system [Who is in charge? Coordination and governance in agentic
software engineering Paolo Ciancarini]

Claims may need to be checked against artifacts, actions
against permissions, and decisions against explicitly assigned authority. [Who is in charge? Coordination and governance in agentic
software engineering Paolo Ciancarini]

Natural language may be vague, incomplete, or inconsistent
Free-form interaction can amplify errors and cause system drift
Shared artifacts, logs, tests, and approval points make decisions traceable [Who is in charge? Coordination and governance in agentic
software engineering Paolo Ciancarini]


Why conversation alone is not enough
Table 2 summarizes the complementary strengths and limitations of conversation as a coordination medium.
The core risk is mistaking fluency for reliability. Because an agent cannot be assumed to distinguish verified knowledge from
plausible interpolation [13], and because its verbal explanations may not faithfully reflect its actual reasoning [14], a coordinator
that relies on conversational output alone has no sufficient mechanism for detecting misalignment between an agent’s stated intent
and its actual computation. [Who is in charge? Coordination and governance in agentic
software engineering Paolo Ciancarini]

-------------------------------------------

[Bo2025 To Rely or Not to Rely?]

LLMs are
prone to hallucination and error, yet often present their outputs
with impressive detail and confidence [13, 90]

Research
in the tangential field of human-AI collaboration has explored
methods like explanations and calibration techniques to improve
joint decision-making by enhancing AI transparency and verifia-
bility, though results often show mixed outcomes and trade-offs
[4, 48, 58, 68, 78]







